% !TEX program = pdflatex
% Example: Book using book class
% Compile with: pdflatex main.tex && bibtex main && pdflatex main.tex && pdflatex main.tex

\documentclass[10pt, oneside, a4paper]{book}

\usepackage{graphicx}
\usepackage{amsmath,amssymb}
\usepackage[margin=2.5cm]{geometry}
\usepackage{makeidx}
\makeindex
\usepackage{hyperref}

\begin{document}

\frontmatter
\title{Introduction to Machine Learning}
\author{Author Name}
\date{\today}
\maketitle

\tableofcontents

\mainmatter

\chapter{Foundations}
\section{What is Machine Learning?}
Machine learning is a subset of artificial intelligence
that enables systems to learn from data.

The concept of \index{machine learning} machine learning
encompasses \index{supervised learning} supervised learning
and \index{unsupervised learning} unsupervised learning.

\chapter{Supervised Learning}
\section{Linear Regression}
Linear regression models the relationship between
a dependent variable and one or more independent variables:
\begin{equation}
y = \beta_0 + \beta_1 x + \epsilon
\end{equation}

\chapter{Unsupervised Learning}
\section{Clustering}
Clustering groups similar data points together.
We discuss \index{clustering!k-means} k-means and
\index{clustering!hierarchical} hierarchical clustering.

\backmatter
\bibliographystyle{plain}
\bibliography{references}

\printindex

\end{document}
